\documentclass[reqno]{amsart} 
\usepackage{amssymb,latexsym,amsmath,amscd,graphicx,setspace,amsthm,verbatim,enumitem}
\usepackage[margin = 3 cm]{geometry}

\def\upint{\mathchoice%
    {\mkern13mu\overline{\vphantom{\intop}\mkern7mu}\mkern-20mu}%
    {\mkern7mu\overline{\vphantom{\intop}\mkern7mu}\mkern-14mu}%
    {\mkern7mu\overline{\vphantom{\intop}\mkern7mu}\mkern-14mu}%
    {\mkern7mu\overline{\vphantom{\intop}\mkern7mu}\mkern-14mu}%
  \int}
\def\lowint{\mkern3mu\underline{\vphantom{\intop}\mkern7mu}\mkern-10mu\int}


\theoremstyle{plain}
\newtheorem{theorem}{Theorem}[section]
\newtheorem{proposition}{Proposition}
\newtheorem{corollary}{Corollary}
\newtheorem{lemma}{Lemma}
\newtheorem{conjecture}{Conjecture}
\newtheorem{question}{Question}
\newtheorem{problem}{Problem}
      
\theoremstyle{definition}
\newtheorem{definition}{Definition}

\newenvironment{solution1}{\paragraph{\emph{Solution $1$}.}}{\hfill$\square$}
\newenvironment{solution2}{\paragraph{\emph{Solution $2$}.}}{\hfill$\square$}
\newenvironment{solution3}{\paragraph{\emph{Solution $3$}.}}{\hfill$\square$}

\begin{document} 

\title[Homework 13]{Homework 13}

\date{\today} 
\maketitle 




\begin{problem}
For all integers $n,i \ge 1$, let $a_{n,i} \in \mathbb{R}$.  Assume that:
\begin{enumerate}
\item For all $n \ge 1$, one has
$$\lim_{n \to \infty} a_{n,i} = a_{i}. $$
\item There exists a sequence $(b_{i})_{i=1}^{\infty}$ of non-negative real numbers such that 
$$|a_{n,i}| \le b_{i}$$
for all $n,i \ge 1$ and for which 
$$\sum_{i=1}^{\infty}b_{i} < \infty. $$
\end{enumerate}
Show that
$$\lim_{n \to \infty}\sum_{i=1}^{\infty}a_{n,i} = \sum_{i=1}^{\infty}a_{i}. $$
(Hint:  Apply the DCT.)
\end{problem}
\begin{solution1}

\end{solution1}

\begin{problem}
Let us revisit the classical equality
$$\lim_{m \to \infty}\left(1 + \frac{x}{m} \right)^{m} = \sum_{n=0}^{\infty}\frac{x^{n}}{n!} $$
valid for all $x \in \mathbb{R}$.  (First, remember for yourself why the series above on the right is absolutely convergent for all $x \in \mathbb{R}$...)
\begin{enumerate}
\item For a fixed $x \in \mathbb{R}$, use the binomial theorem to show that
$$\left(1 + \frac{x}{m} \right)^{m} = \sum_{n=0}^{m} \binom{m}{n} \frac{x^{n}}{m^{n}}. $$
\item For a fixed $x \in \mathbb{R}$, let
$$a_{m,n} = \begin{cases} \binom{m}{n}\frac{x^{n}}{m^{n}}, &\text{ if } n \le m;\\
0, &\text{ otherwise.}
 \end{cases} $$
\item For a fixed $n \in \mathbb{N}$, show that
$$\lim_{m \to \infty} a_{m,n} = \frac{x^{n}}{n!}.$$
\item Show that for all $m,n \ge 1$, one has
$$|a_{m,n}| \le \frac{|x|^{n}}{n!}. $$
\item Using Problem $1$, show that
$$\lim_{m \to \infty}\left(1 + \frac{x}{m} \right)^{m} = \sum_{n=0}^{\infty}\frac{x^{n}}{n!}. $$
\end{enumerate}
\end{problem}
\begin{solution1}

\end{solution1}


\begin{problem}
Let $(X,\mathcal{A},\mu)$ be a measure space and let $f \in \mathcal{L}^{1}(X,\mu)$.  For $i \in \mathbb{N}$, let $A_{i} \in \mathcal{A}$ be such that
$$A_{1} \subseteq A_{2} \subseteq \ldots \subseteq A_{n} \subseteq \ldots $$
and let
$$A = \bigcup_{i=1}^{\infty}A_{i}. $$
Show that
$$\lim_{n \to \infty} \int_{A_{n}}f = \int_{A}f. $$
(Hint:  Consider $f_{n} = f \chi_{A_{n}}$ and apply DCT.)
\end{problem}
\begin{solution1}

\end{solution1}



\begin{problem}
For $i \in \mathbb{N}$, let $a_{i} \in \mathbb{R}$ and assume that
$$\sum_{i=1}^{\infty}a_{i} $$
converges absolutely.  Let $\sigma:\mathbb{N} \rightarrow \mathbb{N}$ be any permutation (i.e. a bijection).  Show that
$$\sum_{i=1}^{\infty}a_{i} = \sum_{i=1}^{\infty}a_{\sigma(i)}. $$
(Hint:  For each $i \in \mathbb{N}$, consider $A_{i} = \{ \sigma(1),\ldots,\sigma(i)\}$ and apply the previous problem.)
\end{problem}
\begin{solution1}

\end{solution1}

This last problem is far from being true if the series is not assumed to be absolutely convergent.  Google {\bfseries Riemann's rearrangement theorem} on Friday night...

In class, we went over improper Riemann integrals carefully and connected this with the Lebesgue theory.  Make sure to do all of the ``optional'' homework problems that I suggested in class... In particular, we showed that we obtain a function $\Gamma:(0,\infty) \rightarrow [0,\infty)$ given by
$$s \mapsto \Gamma(s) = \int_{0}^{\infty} e^{-x}x^{s-1}, $$
and that the integral on the right can be thought of as an improper Riemann integral or as the Lebesgue integral
$$\int_{(0,\infty)} e^{-x}x^{s-1}.$$
Below, you will prove elementary properties of the gamma function.

\begin{problem}
\hfill
\begin{enumerate}
\item Show that $\Gamma(1) = 1$.
\item Using integration by parts, show that 
$$\Gamma(s+1) = s\Gamma(s) $$
for all $s > 0$.
\item Show that for all $n \in \mathbb{N}$, one has
$$\Gamma(n+1) = n!, $$
so that the function $\Gamma$ interpolates the factorial function to non-integer real numbers.
\end{enumerate}
\end{problem}
\begin{solution1}

\end{solution1}

In your free time, think about how one can use the functional equation of the Gamma function to extend its definition to all of $\mathbb{R} \smallsetminus \{\ldots,-3, -2, -1, 0 \}$.  After the break, we will try to understand if $\Gamma$ is a continuous or even differentiable function using the Lebesgue theory.  The same ideas will apply to other integral transforms such as the Fourier, the Laplace and the Mellin transforms.


\end{document}









